\documentclass[
  spanish,
  aspectratio=169,
  xcolor={dvipsnames,table}
]{beamer}

\geometry{paperwidth=19.2cm,paperheight=10.8cm}
\usepackage[utf8]{inputenc}
\usepackage[T1]{fontenc}
\usepackage[spanish,es-tabla]{babel}
\usepackage[scaled=.96]{helvet}
\renewcommand{\familydefault}{\sfdefault}
\usepackage{booktabs}
\usepackage{graphicx}
\usepackage{ragged2e}
\usepackage{tikz}
\usepackage{hyperref}
\usetikzlibrary{calc,positioning,fit,shadows.blur}

\newcommand{\studentname}{Martin Jonathan de la Cruz}
\newcommand{\studentshort}{M. J. de la Cruz}
\newcommand{\studentid}{ES2611202040}
\newcommand{\universityname}{Universidad Abierta y a Distancia de Mexico}
\newcommand{\facultyname}{Licenciatura en Derecho}
\newcommand{\coursename}{Interaprendizajes en ambientes virtuales}
\newcommand{\coursecode}{LDE-IAV-S2B1}
\newcommand{\activitytitle}{Interaprendizajes en ambientes virtuales}
\newcommand{\activitysubtitle}{Encuadre, evaluacion y proyecto colaborativo}
\newcommand{\activityweek}{Panorama de la asignatura}
\newcommand{\teachingfigure}{Janeth Santana Ramirez}
\newcommand{\deliverydate}{\today}
\newcommand{\locationname}{Roma Norte, Ciudad de Mexico}
\newcommand{\departmentlogo}{img/departamentos/UnADM.pdf}

\definecolor{unadmGreenDark}{HTML}{174A3A}
\definecolor{unadmGreen}{HTML}{5F8F3A}
\definecolor{unadmGold}{HTML}{B88A2A}
\definecolor{unadmPaper}{HTML}{F6F7F2}
\definecolor{unadmInk}{HTML}{1F2A24}

\mode<presentation>{
  \usetheme{default}
  \usefonttheme{professionalfonts}
  \setbeamertemplate{navigation symbols}{}
  \setbeamertemplate{blocks}[rounded][shadow=false]
}
\setbeamersize{text margin left=0.60cm,text margin right=0.60cm}
\setbeamercolor{background canvas}{bg=white}
\setbeamercolor{normal text}{fg=unadmInk,bg=white}
\setbeamercolor{structure}{fg=unadmGreenDark}
\setbeamercolor{frametitle}{fg=white,bg=unadmGreenDark}
\setbeamercolor{block title}{fg=white,bg=unadmGreen}
\setbeamercolor{block body}{fg=unadmInk,bg=unadmPaper}
\setbeamerfont{title}{size=\LARGE,series=\bfseries}
\setbeamerfont{frametitle}{size=\large,series=\bfseries}
\setbeamertemplate{itemize item}{\textcolor{unadmGreen}{\large$\blacktriangleright$}}

\setbeamertemplate{frametitle}{
  \nointerlineskip
  \begin{beamercolorbox}[wd=\textwidth,ht=1.02cm,dp=0.22cm,leftskip=0.35cm,rightskip=0.25cm]{frametitle}
    \insertframetitle\hfill\includegraphics[height=0.60cm]{\departmentlogo}
  \end{beamercolorbox}
  {\color{unadmGold}\rule{\textwidth}{1.3pt}}
}

\setbeamertemplate{footline}{
  \leavevmode
  \hbox{
    \begin{beamercolorbox}[wd=.34\paperwidth,ht=2.7ex,dp=1ex,leftskip=1em]{author in head/foot}\color{white}\studentshort\end{beamercolorbox}
    \begin{beamercolorbox}[wd=.40\paperwidth,ht=2.7ex,dp=1ex,center]{title in head/foot}\color{white}\coursename\end{beamercolorbox}
    \begin{beamercolorbox}[wd=.26\paperwidth,ht=2.7ex,dp=1ex,rightskip=1em plus 1fil]{date in head/foot}\color{white}\coursecode\hfill\insertframenumber/\inserttotalframenumber\end{beamercolorbox}
  }
}
\setbeamercolor{author in head/foot}{bg=unadmGreenDark}
\setbeamercolor{title in head/foot}{bg=unadmGreen}
\setbeamercolor{date in head/foot}{bg=unadmGreenDark}

\title[\coursecode]{\activitytitle}
\subtitle{\activitysubtitle}
\author[\studentshort]{\studentname\\\footnotesize Matricula: \studentid}
\institute[UnADM]{\universityname\\\facultyname}
\date[\activityweek]{\locationname\\\activityweek\\\deliverydate}

\begin{document}

\begin{frame}[plain]
  \begin{tikzpicture}[remember picture,overlay]
    \fill[unadmGreenDark] (current page.south west) rectangle ([xshift=0.58\paperwidth]current page.north west);
    \fill[unadmGreen] ([xshift=0.58\paperwidth]current page.south west) rectangle (current page.north east);
    \node[anchor=center,opacity=0.10] at (current page.center) {\includegraphics[width=0.72\paperwidth]{\departmentlogo}};
    \node[anchor=north west] at ([xshift=0.58cm,yshift=-0.46cm]current page.north west) {\includegraphics[width=2.55cm]{\departmentlogo}};
  \end{tikzpicture}
  \vspace*{1.80cm}
  {\color{white}\fontsize{25}{30}\selectfont\bfseries \activitytitle\par}
  \vspace{0.28cm}
  {\color{white!86}\large \activitysubtitle}\\[0.35cm]
  {\color{unadmGold}\rule{0.58\linewidth}{1.3pt}}\\[0.35cm]
  {\color{white}\studentname}\\
  {\color{white!82}\footnotesize \facultyname}
\end{frame}

\begin{frame}{Objetivo y alcance}
  \begin{block}{Objetivo editable}
    Construir estrategias de interaprendizaje virtual para colaborar,
    comunicarse con respeto, elegir TICCAD utiles y producir trabajo juridico
    verificable en entornos digitales.
  \end{block}
  \begin{itemize}
    \item Optativa de segundo semestre, bloque 1, en la Licenciatura en Derecho.
    \item La participacion se entiende como dialogo, colaboracion y construccion conjunta.
    \item La figura docente acompana el proceso y promueve autonomia y respeto.
  \end{itemize}
\end{frame}

\begin{frame}{Ejes de la materia}
  \begin{itemize}
    \item Diversidad cultural y linguistica como marco del trabajo en linea.
    \item Interaprendizaje para distinguir colaboracion real de reparto mecanico.
    \item TICCAD como ecosistema minimo, accesible y trazable.
    \item Bitacora, autoevaluacion y coevaluacion con evidencia concreta.
  \end{itemize}
\end{frame}

\begin{frame}{Ruta de evaluacion}
  \begin{block}{Secuencia con mayor peso}
    El curso prepara un despacho juridico digital: primero sensibiliza y ordena
    la colaboracion; despues consolida un producto profesional colectivo.
  \end{block}
  \begin{itemize}
    \item Foro de diversidad cultural y linguistica: 10 puntos.
    \item Interaprendizaje: comprender, construir y distinguir: 10 puntos.
    \item Ecosistema TICCAD: 15 puntos.
    \item Documento colaborativo del despacho juridico digital: 20 puntos.
    \item Propuesta profesional del despacho juridico digital: 30 puntos.
    \item Bitacora de interaprendizaje: 15 puntos.
  \end{itemize}
\end{frame}

\begin{frame}{Proyecto integrador}
  \begin{block}{Despacho juridico digital}
    La materia culmina en una simulacion profesional donde el equipo organiza
    roles, acuerdos, herramientas y una propuesta vinculada con un caso.
  \end{block}
  \begin{itemize}
    \item La colaboracion debe dejar rastro verificable en acuerdos y entregables.
    \item Las TICCAD se justifican por funcion, accesibilidad y control de versiones.
    \item La coevaluacion solo vale si se apoya en hechos observables.
  \end{itemize}
\end{frame}

\begin{frame}{Fuentes y lectura editorial}
  \begin{itemize}
    \item Bibliografia local: \texttt{optativa-semestre-2-bloque-1.bib}.
    \item Fuente curricular: malla de Derecho UnADM en assets-unadm.
    \item Insumos base: notas de bienvenida e introduccion materializadas en esta carpeta.
    \item Revisar siempre coherencia entre participacion, evidencia y conclusion.
  \end{itemize}
\end{frame}

\end{document}